\documentclass[../../main.tex]{subfiles}% 注意这里的文件路径不能用 ./main.tex ,否则用latexmk编译子文件会报错
\graphicspath{{\subfix{./image/}}{\subfix{../../image/}}} % 指定图片引用路径,后续可以直接使用图片文件名
% 如果图片存放在上级目录,则使用 ../../image/ 作为备用路径

% 例如:
% \begin{figure}[H]
% \centering
% \includegraphics[width=0.3\textwidth]{图.png}
% \caption{}
% \label{figure:图}
% \end{figure}
% 注意:上述\label{}一定要放在\caption{}之后,否则引用图片序号会只会显示??.

\begin{document}

\section{初值问题}

\subsection{Fourier 变换和 Fourier 积分}

\begin{definition}
设 $f(x)\in L^1(\mathbb{R})$，则无穷积分
\begin{align}
\hat{f}(\lambda)=\frac{1}{\sqrt{2\pi}}\int_{-\infty}^{+\infty}f(x)e^{-\mathrm{i}\lambda  x}\,\mathrm{d}x \label{eq:fourier-def}
\end{align}
有意义，称为 $f(x)$ 的 \textbf{Fourier 变换}，记为 $(f(x))^\wedge=\hat{f}(\lambda)$。
\end{definition}

\begin{theorem}[Fourier积分定理]
设 $f(x)\in L^1(\mathbb{R})\cap C^1(\mathbb{R})$，则
\begin{align}
f(x)=\lim_{N\rightarrow +\infty} \frac{1}{\sqrt{2\pi}}\int_{-N}^N{\hat{f}(\lambda )e^{\mathrm{i}\lambda x}\,\mathrm{d}\lambda}=\frac{1}{\sqrt{2\pi}}\int_{-\infty}^{+\infty}{\hat{f}(\lambda )e^{\mathrm{i}\lambda x}\,\mathrm{d}\lambda}.\label{eq:fourier-inversion}
\end{align}
公式\eqref{eq:fourier-inversion}称为\textbf{反演公式},右端的积分表示在 Cauchy 主值意义下的无穷积分，通常称为 \textbf{Fourier 逆变换}，记为 $(\hat{f}(\lambda))^\vee$。因此 Fourier 积分定理也可写成
\begin{align*}
(\hat{f})^\vee=f.
\end{align*}
\end{theorem}
\begin{remark}
也就是说，$L^1(\mathbb{R})\cap C^1(\mathbb{R})$ 中的函数作一次 Fourier 变换后，再作一次 Fourier 逆变换，又回到函数本身.
\end{remark}
\begin{proof}
由于 $f(x)\in L^1(\mathbb{R})$，则对任意的 $\lambda\in\mathbb{R}$，无穷积分
\begin{align*}
\int_{-\infty}^{+\infty}f(x)e^{-\mathrm{i}\lambda  x}\,\mathrm{d}x
\end{align*}
一致收敛，且是 $\lambda$ 的连续函数。固定 $x\in\mathbb{R}$，我们利用上述无穷积分的一致收敛性，交换积分次序得
\begin{align*}
\frac{1}{\sqrt{2\pi}}\int_{-N}^{N}\hat{f}(\lambda)e^{\mathrm{i}\lambda  x}\,\mathrm{d}\lambda
&=\frac{1}{2\pi}\int_{-\infty}^{+\infty}f(\xi)\,\mathrm{d}\xi\int_{-N}^{N}e^{\mathrm{i}\lambda (x-\xi)}\,\mathrm{d}\lambda \\
&=\frac{1}{\pi}\int_{-\infty}^{+\infty}f(\xi)\frac{\sin N(x-\xi)}{x-\xi}\,\mathrm{d}\xi \\
&=\frac{1}{\pi}\int_{-\infty}^{+\infty}f(x+\eta)\frac{\sin N\eta}{\eta}\,\mathrm{d}\eta.
\end{align*}
设 $M$ 是一个待定正数。我们记
\begin{align*}
I_1=\frac{1}{\pi}\int_{-\infty}^{-M}f(x+\eta)\frac{\sin N\eta}{\eta}\,\mathrm{d}\eta,\\
I_2=\frac{1}{\pi}\int_{-M}^{M}f(x+\eta)\frac{\sin N\eta}{\eta}\,\mathrm{d}\eta,\\
I_3=\frac{1}{\pi}\int_{M}^{+\infty}f(x+\eta)\frac{\sin N\eta}{\eta}\,\mathrm{d}\eta,
\end{align*}
则
\begin{align*}
\frac{1}{\sqrt{2\pi}}\int_{-N}^{N}\hat{f}(\lambda)e^{\mathrm{i}\lambda  x}\,\mathrm{d}\lambda=I_1+I_2+I_3.
\end{align*}
由正弦函数的有界性我们估计 $I_1$ 和 $I_3$ 得到
\begin{align*}
|I_1|+|I_3|\leqslant \frac{1}{\pi M}\int_{-\infty}^{+\infty}|f(\eta)|\,\mathrm{d}\eta.
\end{align*}
接着我们估计 $I_2$：
\begin{align*}
I_2&=\frac{1}{\pi}\int_{-M}^{M}\frac{f(x+\eta)-f(x)}{\eta}\sin N\eta\,\mathrm{d}\eta+\frac{f(x)}{\pi}\int_{-M}^{M}\frac{\sin N\eta}{\eta}\,\mathrm{d}\eta \\
&=\frac{1}{\pi}\int_{-M}^{M}g(x,\eta)\sin N\eta\,\mathrm{d}\eta+\frac{f(x)}{\pi}\int_{-MN}^{MN}\frac{\sin\eta}{\eta}\,\mathrm{d}\eta,
\end{align*}
其中
\begin{align*}
g(x,\eta)=\int_0^1 f'(x+\tau\eta)\,\mathrm{d}\tau
\end{align*}
是关于 $\eta$ 的连续函数。对任意的 $\varepsilon>0$，先取定正数 $M$，使得
\begin{align*}
|I_1|+|I_3|<\frac{\varepsilon}{2},
\end{align*}
然后利用\hyperref[Basis of Analytics-lemma:Riemann-Lebesgue引理]{Riemann-Lebesgue引理}和\rrefpro{Basis of Analytics-proposition:重要定积分结论}{Basis of Analytics-proposition:重要定积分结论-2}再取定 $N_0$，使得当 $N>N_0$ 时，
\begin{align*}
\left|\frac{1}{\pi}\int_{-M}^{M}g(x,\eta)\sin N\eta\,\mathrm{d}\eta\right|<\frac{\varepsilon}{4},
\end{align*}
和
\begin{align*}
\left|\frac{f(x)}{\pi}\int_{-MN}^{MN}\frac{\sin\eta}{\eta}\,\mathrm{d}\eta-f(x)\right|<\frac{\varepsilon}{4}.
\end{align*}
综上所述，我们得到，当 $N>N_0$ 时，
\begin{align*}
\left|\frac{1}{\sqrt{2\pi}}\int_{-N}^{N}\hat{f}(\lambda)e^{\mathrm{i}\lambda  x}\,\mathrm{d}\lambda-f(x)\right|\leqslant \varepsilon.
\end{align*}
从而完成定理的证明.
\end{proof}

\begin{proposition}[$\mathbb{R}$上的Fourier 变换的基本性质]\label{proposition:Fourier 变换的基本性质1}
设 $f,g\in L^1(\mathbb{R})$，$a_1,a_2\in\mathbb{C}$，$k\in\mathbb{R}$，$k\neq 0$，则其 Fourier 变换
\begin{align*}
\hat{f}(\lambda)\triangleq \frac{1}{\sqrt{2\pi}}\int_{-\infty}^{+\infty}f(x)e^{-\mathrm{i}\lambda x}\,\mathrm{d}x
\end{align*}
满足以下性质：
\begin{enumerate}[(1)]
\item\label{proposition:Fourier 变换的基本性质1-1} \textbf{线性性质:}
\begin{align*}
(a_1f_1+a_2f_2)^\wedge=a_1\hat{f}_1+a_2\hat{f}_2.
\end{align*}

\item\label{proposition:Fourier 变换的基本性质1-2} \textbf{微商性质:}若 $f,f'\in L^1(\mathbb{R})\cap C(\mathbb{R})$，则
\begin{align*}
(f')^\wedge=\mathrm{i}\lambda \hat{f}(\lambda).
\end{align*}
更一般地，若 $f,f',\cdots,f^{(m)}\in L^1(\mathbb{R})\cap C(\mathbb{R})$，则
\begin{align*}
(f^{(m)})^\wedge=(\mathrm{i}\lambda)^m\hat{f}(\lambda).
\end{align*}

\item\label{proposition:Fourier 变换的基本性质1-3} \textbf{乘多项式性质:}若 $f,xf\in L^1(\mathbb{R})$，则
\begin{align*}
(xf)^\wedge=\mathrm{i}\frac{\mathrm{d}}{\mathrm{d}\lambda}\hat{f}(\lambda).
\end{align*}
更一般地，若 $f,xf,\cdots,x^m f\in L^1(\mathbb{R})$，则
\begin{align*}
(x^m f)^\wedge=\mathrm{i}^m\frac{\mathrm{d}^m}{\mathrm{d}\lambda^m}\hat{f}(\lambda).
\end{align*}

\item\label{proposition:Fourier 变换的基本性质1-4} \textbf{平移性质:}
\begin{align*}
(f(x-k))^\wedge=e^{-\mathrm{i}\lambda k}\hat{f}(\lambda).
\end{align*}

\item\label{proposition:Fourier 变换的基本性质1-5} \textbf{伸缩性质:}
\begin{align*}
(f(kx))^\wedge=\frac{1}{|k|}\hat{f}\left(\frac{\lambda}{k}\right).
\end{align*}

\item\label{proposition:Fourier 变换的基本性质1-6} \textbf{对称性质:}
\begin{align*}
(f)^\vee=\hat{f}(-\lambda),
\end{align*}
其中 $(f)^\vee$ 表示 Fourier 逆变换.

\item\label{proposition:Fourier 变换的基本性质1-7} \textbf{卷积性质:}若$f,g\in L^1(\mathbb{R})$,定义卷积
\begin{align*}
(f*g)(x)\triangleq\int_{-\infty}^{+\infty}f(x-t)g(t)\,\mathrm{d}t,
\end{align*}
则 $f*g\in L^1(\mathbb{R})$，且
\begin{align*}
(f*g)^\wedge=\sqrt{2\pi}\,\hat{f}(\lambda)\hat{g}(\lambda).
\end{align*}
\end{enumerate}
\end{proposition}
\begin{proof}
\begin{enumerate}[(1)]
\item 

\item 由 $f'\in C(\mathbb{R})\cap L^1(\mathbb{R})$ 及 Newton-Leibniz 公式，极限 $\lim_{|x|\to\infty}f(x)=0$ 存在且为零。于是分部积分得
\begin{align*}
(f')^\wedge
=\frac{1}{\sqrt{2\pi}}\int_{-\infty}^{+\infty}f'(x)e^{-\mathrm{i}\lambda x}\,\mathrm{d}x 
=\frac{1}{\sqrt{2\pi}}\left[f(x)e^{-\mathrm{i}\lambda x}\right]_{-\infty}^{+\infty}
+\frac{\mathrm{i}\lambda}{\sqrt{2\pi}}\int_{-\infty}^{+\infty}f(x)e^{-\mathrm{i}\lambda x}\,\mathrm{d}x 
=\mathrm{i}\lambda \hat{f}(\lambda).
\end{align*} 

\item  由于 $f,xf\in L^1(\mathbb{R})$，$\hat{f}$ 可微且
\begin{align*}
\frac{\mathrm{d}}{\mathrm{d}\lambda}\hat{f}(\lambda)
&=\frac{1}{\sqrt{2\pi}}\int_{-\infty}^{+\infty}f(x)(-\mathrm{i}x)e^{-\mathrm{i}\lambda x}\,\mathrm{d}x
=-\mathrm{i}(xf)^\wedge,
\end{align*}
故 $(xf)^\wedge=\mathrm{i}\frac{\mathrm{d}}{\mathrm{d}\lambda}\hat{f}(\lambda)$.

\item 注意到
\begin{align*}
(f(x-k))^\wedge
=\frac{1}{\sqrt{2\pi}}\int_{-\infty}^{+\infty}f(x-k)e^{-\mathrm{i}\lambda x}\,\mathrm{d}x 
=\frac{1}{\sqrt{2\pi}}\int_{-\infty}^{+\infty}f(y)e^{-\mathrm{i}\lambda(y+k)}\,\mathrm{d}y
=e^{-\mathrm{i}\lambda k}\hat{f}(\lambda).
\end{align*}

\item  以 $k>0$ 为例，作变量替换 $y=kx$：
\begin{align*}
(f(kx))^\wedge
=\frac{1}{\sqrt{2\pi}}\int_{-\infty}^{+\infty}f(kx)e^{-\mathrm{i}\lambda x}\,\mathrm{d}x 
=\frac{1}{\sqrt{2\pi}}\int_{-\infty}^{+\infty}f(y)e^{-\mathrm{i}\lambda y/k}\,\frac{\mathrm{d}y}{k}
=\frac{1}{k}\hat{f}\left(\frac{\lambda}{k}\right).
\end{align*}
$k<0$ 时同理得 $\frac{1}{|k|}\hat{f}(\lambda/k)$.

\item  

\item 首先
\begin{align*}
\int_{-\infty}^{+\infty}|f*g(x)|\,\mathrm{d}x
\leqslant \int_{-\infty}^{+\infty}\int_{-\infty}^{+\infty}|f(x-t)g(t)|\,\mathrm{d}t\,\mathrm{d}x 
=\int_{-\infty}^{+\infty}|g(t)|\,\mathrm{d}t\int_{-\infty}^{+\infty}|f(x-t)|\,\mathrm{d}x
=\|f\|_1\|g\|_1,
\end{align*}
故 $f*g\in L^1(\mathbb{R})$. 再由\hyperref[Real Analysis-theorem:Fubini定理 可积函数的情形]{Fubini定理}得
\begin{align*}
(f*g)^\wedge
&=\frac{1}{\sqrt{2\pi}}\int_{-\infty}^{+\infty}e^{-\mathrm{i}\lambda x}\int_{-\infty}^{+\infty}f(x-t)g(t)\,\mathrm{d}t\,\mathrm{d}x \\
&=\frac{1}{\sqrt{2\pi}}\int_{-\infty}^{+\infty}g(t)e^{-\mathrm{i}\lambda t}
\left(\int_{-\infty}^{+\infty}f(x-t)e^{-\mathrm{i}\lambda (x-t)}\,\mathrm{d}x\right)\mathrm{d}t \\
&=\sqrt{2\pi}\,\hat{f}(\lambda)\hat{g}(\lambda).
\end{align*}
\end{enumerate}
\end{proof}

\begin{example}
设
\begin{align*}
f_1(x)=\begin{cases}
1, & |x|\leqslant A,\\
0, & |x|>A
\end{cases}\quad (A>0).
\end{align*}
求 $\hat{f}_1(\lambda)$。
\end{example}
\begin{remark}
此例可以用于一维波动方程初值问题的求解。
\end{remark}
\begin{solution}
由定义，有
\begin{align*}
\hat{f}_1(\lambda)=\frac{1}{\sqrt{2\pi}}\int_{-A}^{A}e^{-\mathrm{i}\lambda  x}\,\mathrm{d}x=\sqrt{\frac{2}{\pi}}\frac{\sin\lambda A}{\lambda}.
\end{align*}
\end{solution}

\begin{example}
设
\begin{align*}
f_2(x)=\begin{cases}
e^{-x}, & x\geqslant 0,\\
0, & x<0.
\end{cases}
\end{align*}
求 $\hat{f}_2(\lambda)$。
\end{example}
\begin{solution}
由定义，有
\begin{align*}
\hat{f}_2(\lambda)=\frac{1}{\sqrt{2\pi}}\int_0^{+\infty}e^{-(1+\mathrm{i}\lambda )x}\,\mathrm{d}x=\frac{1}{\sqrt{2\pi}(1+\mathrm{i}\lambda )}.
\end{align*}
\end{solution}

\begin{example}
设 $f_3(x)=e^{-|x|}$。求 $\hat{f}_3(\lambda)$。
\end{example}
\begin{solution}
由于 $f_3(x)=f_2(x)+f_2(-x)$，利用\hyperref[proposition:Fourier 变换的基本性质1-1]{线性性质}和\hyperref[proposition:Fourier 变换的基本性质1-5]{伸缩性质}则有
\begin{align*}
\hat{f}_3(\lambda)=\hat{f}_2(\lambda)+(f_2(-x))^\wedge 
=\hat{f}_2(\lambda)+\hat{f}_2(-\lambda) 
=\frac{1}{\sqrt{2\pi}(1+\mathrm{i}\lambda )}+\frac{1}{\sqrt{2\pi}(1-\mathrm{i}\lambda )} 
=\frac{2}{\sqrt{2\pi}(1+\lambda^2)}.
\end{align*}
\end{solution}

\begin{example}
设 $f_4(x)=e^{-x^2}$。求 $\hat{f}_4(\lambda)$。
\end{example}
\begin{solution}
由分部积分公式和\hyperref[proposition:Fourier 变换的基本性质1-3]{乘多项式性质}我们得到，当 $\lambda\neq 0$ 时，
\begin{align*}
\hat{f}_4(\lambda)&=\frac{1}{\sqrt{2\pi}}\int_{-\infty}^{+\infty}e^{-x^2}e^{-\mathrm{i}\lambda  x}\,\mathrm{d}x \\
&=-\frac{1}{\sqrt{2\pi}}\cdot\frac{1}{\mathrm{i}\lambda }\left(e^{-x^2}e^{-\mathrm{i}\lambda  x}\Big|_{-\infty}^{+\infty}+2\int_{-\infty}^{+\infty}xe^{-x^2}e^{-\mathrm{i}\lambda  x}\,\mathrm{d}x\right) \\
&=\frac{2i}{\lambda}(xf_4(x))^\wedge=-\frac{2}{\lambda}\cdot\frac{\mathrm{d}}{\mathrm{d}\lambda}\hat{f}_4(\lambda).
\end{align*}
注意到恒等式
\begin{align*}
\int_{-\infty}^{+\infty}e^{-x^2}\,\mathrm{d}x=\sqrt{\pi},
\end{align*}
我们导出 $\hat{f}_4(\lambda)$ 作为 $\lambda$ 的函数满足下列常微分方程初值问题
\begin{align*}
\begin{cases}
\frac{\mathrm{d}}{\mathrm{d}\lambda}\hat{f}_4(\lambda)=-\frac{\lambda}{2}\hat{f}_4(\lambda),\\
\hat{f}_4(0)=\frac{1}{\sqrt{2}}.
\end{cases}
\end{align*}
求解上述常微分方程初值问题得 $\hat{f}_4(\lambda)=\frac{1}{\sqrt{2}}e^{-\lambda^2/4}$。
\end{solution}

\begin{example}
设
\begin{align*}
f_5(x)=e^{-Ax^2}\quad (A>0).
\end{align*}
求 $\hat{f}_5(\lambda)$。
\end{example}
\begin{solution}
利用\hyperref[proposition:Fourier 变换的基本性质1-5]{伸缩性质}得到
\begin{align*}
\hat{f}_5(\lambda)=(f_4(\sqrt{A}x))^\wedge=\frac{1}{\sqrt{A}}\hat{f}_4\left(\frac{\lambda}{\sqrt{A}}\right)
=\frac{1}{\sqrt{2A}}e^{-\frac{\lambda^2}{4A}}.
\end{align*}
\end{solution}

\begin{definition}
设 $f(x)=f(x_1,x_2,\cdots,x_n)\in L^1(\mathbb{R}^n)$，则无穷积分
\begin{align}
\hat{f}(\lambda)=\frac{1}{(\sqrt{2\pi})^n}\int_{\mathbb{R}^n}f(x)e^{-\mathrm{i}\lambda \cdot x}\,\mathrm{d}x \label{eq:fourier-multidim}
\end{align}
有意义，称为 $f(x)$ 的 \textbf{Fourier 变换}，这里 $\lambda=(\lambda_1,\lambda_2,\cdots,\lambda_n)\in\mathbb{R}^n$，记为 $(f(x))^\wedge=\hat{f}(\lambda)$。
\end{definition}

\begin{theorem}[反演公式]
设 $f(x)\in L^1(\mathbb{R}^n)\cap C^1(\mathbb{R}^n)$，则
\begin{align}
(\hat{f}(\lambda))^\vee=\lim_{N\to+\infty}\frac{1}{(\sqrt{2\pi})^n}\int_{|\lambda|\leqslant N}\hat{f}(\lambda)e^{\mathrm{i}\lambda \cdot x}\,\mathrm{d}\lambda=f(x), \label{eq:fourier-inv-multidim}
\end{align}
其中 $x=(x_1,x_2,\cdots,x_n)\in\mathbb{R}^n$。$(\hat{f}(\lambda))^\vee$ 表示在 Cauchy 主值意义下的无穷积分，称为 $\hat{f}(\lambda)$ 的 \textbf{Fourier 逆变换}。
\end{theorem}

\begin{proposition}[$\mathbb{R}^n$上的Fourier变换的基本性质]\label{proposition:Fourier 变换的基本性质2}
设 $f,g\in L^1(\mathbb{R}^n)$，$a_1,a_2\in\mathbb{C}$，$k\in\mathbb{R}\setminus\{0\}$，$x_0\in\mathbb{R}^n$，则其Fourier变换
\begin{align*}
\hat{f}(\lambda)\triangleq \frac{1}{(2\pi)^{n/2}}\int_{\mathbb{R}^n} f(x)e^{-\mathrm{i}\lambda\cdot x}\,\mathrm{d}x
\end{align*}
满足以下性质：
\begin{enumerate}[(1)]
\item \textbf{线性性质:}
\begin{align*}
(a_1f_1+a_2f_2)^\wedge=a_1\hat{f}_1+a_2\hat{f}_2.
\end{align*}

\item \textbf{微商性质:} 若 $f\in C^m(\mathbb{R}^n)$ 且对所有 $|\alpha|\leqslant m$ 有 $D^\alpha f\in L^1(\mathbb{R}^n)$，则
\begin{align*}
(D^\alpha f)^\wedge=(\mathrm{i}\lambda)^\alpha \hat{f}(\lambda),\qquad |\alpha|\leqslant m,
\end{align*}
其中 $\lambda^\alpha=\lambda_1^{\alpha_1}\cdots\lambda_n^{\alpha_n}$.

\item \textbf{乘多项式性质:} 若 $x^\alpha f\in L^1(\mathbb{R}^n)$ 对所有 $|\alpha|\leqslant m$ 成立，则
\begin{align*}
(x^\alpha f)^\wedge=\mathrm{i}^{|\alpha|}D^\alpha \hat{f}(\lambda),\qquad |\alpha|\leqslant m.
\end{align*}

\item \textbf{平移性质:}
\begin{align*}
(f(x-x_0))^\wedge=e^{-\mathrm{i}\lambda\cdot x_0}\hat{f}(\lambda).
\end{align*}

\item \textbf{伸缩性质:}
\begin{align*}
(f(kx))^\wedge=\frac{1}{|k|^n}\hat{f}\left(\frac{\lambda}{k}\right).
\end{align*}

\item \textbf{对称性质:}
\begin{align*}
(f)^\vee=\hat{f}(-\lambda),
\end{align*}
其中 $(f)^\vee$ 表示 Fourier 逆变换.

\item \textbf{卷积性质:} 若$f,g\in L^1(\mathbb{R})$,定义卷积
\begin{align*}
(f*g)(x)\triangleq\int_{\mathbb{R}^n} f(x-y)g(y)\,\mathrm{d}y,
\end{align*}
则 $f*g\in L^1(\mathbb{R}^n)$，且
\begin{align*}
(f*g)^\wedge=(2\pi)^{\frac{n}{2}}\,\hat{f}(\lambda)\hat{g}(\lambda).
\end{align*}

\item\label{proposition:Fourier 变换的基本性质2-8} \textbf{分离变量性质:} 若 $f(x)=\prod_{j=1}^n f_j(x_j)$，其中 $f_j\in L^1(\mathbb{R})$，则
\begin{align*}
\hat{f}(\lambda)=\prod_{j=1}^n \hat{f}_j(\lambda_j).
\end{align*}
\end{enumerate}
\end{proposition}

\begin{example}
设 $f_6(x)=e^{-A|x|^2}$ ($A>0$)。求 $\hat{f}_6(\lambda)$。
\end{example}
\begin{solution}
事实上，由\hyperref[proposition:Fourier 变换的基本性质2-8]{分离变量性质}我们得到
\begin{align*}
\hat{f}_6(\lambda)=\prod_{i=1}^{n}(e^{-Ax_i^2})^\wedge=\prod_{i=1}^{n}\frac{1}{\sqrt{2A}}e^{-\frac{\lambda_i^2}{4A}}=\frac{1}{(2A)^{\frac{n}{2}}}e^{-\frac{|\lambda|^2}{4A}}.
\end{align*}
\end{solution}

\subsection{初值问题和基本解}

\begin{definition}
对 $x\in\mathbb{R}$，$t>0$，称函数
\begin{align*}
K(x,t)\triangleq \begin{cases}
\frac{1}{2a\sqrt{\pi t}}e^{-\frac{x^2}{4a^2t}}, & t>0,\\
0, & t\leqslant 0
\end{cases} 
\end{align*}
为\textbf{一维热方程的Poisson 核}或\textbf{一维热方程的热核}，其中 $a>0$ 为常数.
\end{definition}

\begin{definition}
设 $\varphi\in C(\mathbb{R})$ 有界，$f\in C(\mathbb{R}\times(0,\infty))$ 有界，$a>0$ 为常数. 对任意 $x\in\mathbb{R}$，$t>0$，定义
\begin{align}
u(x,t)\triangleq \int_{-\infty}^{+\infty}K(x-\xi,t)\varphi(\xi)\,\mathrm{d}\xi
+\int_0^t\int_{-\infty}^{+\infty}K(x-\xi,t-\tau)f(\xi,\tau)\,\mathrm{d}\xi\,\mathrm{d}\tau, \label{def:poisson1d}
\end{align}
其中 $K$是一维热方程的Poisson核. 称 \eqref{def:poisson1d}式为\textbf{一维热方程初值问题的Poisson 公式}.
\end{definition}

\begin{definition}
设 $\varGamma $ 是定义在 $\mathbb{R}\times(0,\infty)\times\mathbb{R}\times(0,\infty)$ 上的函数（其中 $t>\tau$），定义为
\begin{align*}
\varGamma (x,t;\xi,\tau)\triangleq K(x-\xi,t-\tau),\quad t>\tau,
\end{align*}
其中 $K$是一维热方程的Poisson核. 即
\begin{align}
\varGamma (x,t;\xi,\tau)=
\begin{cases}
\dfrac{1}{2a\sqrt{\pi(t-\tau)}}\exp\left(-\dfrac{(x-\xi)^2}{4a^2(t-\tau)}\right), & t>\tau,\\[0.4cm]
0, & t\leqslant \tau.
\end{cases}
\end{align}
称 $\varGamma $ 为\textbf{一维热方程的基本解}.
\end{definition}
\begin{remark}
由定义可知，对任意固定的 $(\xi,\tau)$，$\varGamma (x,t;\xi,\tau)$ 作为 $(x,t)$ 的函数在 $t>\tau$ 时满足热方程
\begin{align*}
\partial_t \varGamma  - a^2\partial_x^2\varGamma  = 0,
\end{align*}
且在 $t=\tau$ 处有初始奇性 $\varGamma (x,\tau;\xi,\tau)=\delta(x-\xi)$（在广义函数意义下）.
\end{remark}
\begin{remark}
基本解的物理意义如下：考虑两端均在无穷远点，侧表面绝热的均匀细杆。在 $t=\tau$ 时刻，在 $x=\xi$ 处放置一个瞬时单位点热源，那么由这个瞬时单位点热源在杆上产生的温度分布实际上就是基本解 $\varGamma (x,t;\xi,\tau)$。因此基本解 $\varGamma (x,t;\xi,\tau)$ 也称为\textbf{点源函数}。
\end{remark}
\begin{remark}
热方程的基本解也可如下定义:

对于任意 $(\xi,\tau)\in\mathbb{R}_+^2=\mathbb{R}\times\mathbb{R}_+$，若函数 $u(x,t)\in L_{\mathrm{loc}}^1(\mathbb{R}_+^2)\cap C(\mathbb{R}_+^2\setminus(\xi,\tau))$ 且在广义函数意义下满足方程和初始条件
\begin{align*}
\begin{cases}
\frac{\partial u}{\partial t}-a^2\frac{\partial^2 u}{\partial x^2}=\delta(x-\xi,t-\tau), & (x,t)\in\mathbb{R}\times\mathbb{R}_+,\\
u(x,0)=0, & x\in\mathbb{R},
\end{cases}
\end{align*}
则称之为\textbf{一维热方程的基本解}，且记为 $\varGamma (x,t;\xi,\tau)$。这里 $\delta(x-\xi,t-\tau)$ 是 Dirac 测度。
\end{remark}

\begin{proposition}[热方程基本解的性质]\label{proposition:热方程基本解的性质}
设 $\varGamma(x,t;\xi,\tau)$ 为一维热方程的基本解，则对任意 $x,\xi\in\mathbb{R}$，$t>\tau$，成立以下性质：
\begin{enumerate}[(1)]
\item $\varGamma(x,t;\xi,\tau)>0$.

\item $\varGamma(x,t;\xi,\tau)=\varGamma(\xi,t;x,\tau).$

\item $\int_{-\infty}^{+\infty}\varGamma(x,t;\xi,\tau)\,\mathrm{d}\xi=1.$

\item $\varGamma(x,t;\xi,\tau)$ 关于所有自变量无穷次连续可微，且满足
\begin{align*}
\frac{\partial\varGamma}{\partial t}-a^2\frac{\partial^2\varGamma}{\partial x^2}=0,\qquad 
\frac{\partial\varGamma}{\partial\tau}+a^2\frac{\partial^2\varGamma}{\partial\xi^2}=0.
\end{align*}

\item $|\varGamma(x,t;\xi,\tau)|\leqslant \frac{1}{2a\sqrt{\pi(t-\tau)}}.$

\item 若 $\varphi\in C(\mathbb{R})\cap L^\infty(\mathbb{R})$，则对任意 $x\in\mathbb{R}$，
\begin{align*}
\lim_{t\to 0^+}\int_{-\infty}^{+\infty}\varGamma(x,t;\xi,0)\varphi(\xi)\,\mathrm{d}\xi=\varphi(x).
\end{align*}
\end{enumerate}
\end{proposition}
\begin{proof}
\begin{enumerate}[(1)]
\item 由定义 $\varGamma(x,t;\xi,\tau)=\frac{1}{2a\sqrt{\pi(t-\tau)}}\exp\left(-\frac{(x-\xi)^2}{4a^2(t-\tau)}\right)$，因指数函数恒正且系数为正，故 $\varGamma>0$.

\item 由定义直接验证可得.

\item 作变量替换 $\eta=\frac{\xi-x}{2a\sqrt{t-\tau}}$，则
\begin{align*}
\int_{-\infty}^{+\infty}\varGamma(x,t;\xi,\tau)\,\mathrm{d}\xi
&=\frac{1}{\sqrt{\pi}}\int_{-\infty}^{+\infty}e^{-\eta^2}\,\mathrm{d}\eta=1.
\end{align*}

\item 直接计算可得.

\item 由 $\varGamma$ 的表达式，$e^{-\frac{(x-\xi)^2}{4a^2(t-\tau)}}\leqslant 1$，故
\begin{align*}
|\varGamma(x,t;\xi,\tau)|\leqslant \frac{1}{2a\sqrt{\pi(t-\tau)}}.
\end{align*}

\item 作变量替换 $\eta=\frac{\xi-x}{2a\sqrt{t}}$，则
\begin{align*}
\int_{-\infty}^{+\infty}\varGamma(x,t;\xi,0)\varphi(\xi)\,\mathrm{d}\xi
=\frac{1}{\sqrt{\pi}}\int_{-\infty}^{+\infty}e^{-\eta^2}\varphi(x+2a\sqrt{t}\eta)\,\mathrm{d}\eta.
\end{align*}
由 $\varphi$ 的有界性和连续性，积分关于 $t$ 一致收敛，故可交换极限与积分：
\begin{align*}
\lim_{t\to 0^+}\int_{-\infty}^{+\infty}\varGamma(x,t;\xi,0)\varphi(\xi)\,\mathrm{d}\xi
&=\frac{1}{\sqrt{\pi}}\int_{-\infty}^{+\infty}e^{-\eta^2}\varphi(x)\,\mathrm{d}\eta
=\varphi(x).
\end{align*}
\end{enumerate}
\end{proof}

\begin{theorem}
设 $\varphi(x)\in C(\mathbb{R})\cap L^\infty(\mathbb{R})$，$f(x,t)\equiv 0$，则一维热方程初值问题的Poisson公式\eqref{def:poisson1d}表示的函数
\begin{align*}
u(x,t)=\int_{-\infty}^{+\infty}K(x-\xi,t)\varphi(\xi)\,\mathrm{d}\xi
\end{align*}
是一维热方程初值问题\eqref{eq:heat-cauchy1d} 的有界 $C^{2,1}(\mathbb{R}_+^2)$ 解，其中函数 $u(x,t)$ 在极限的意义下取初值 $\varphi(x)$。
\end{theorem}
\begin{remark}
定理中关于 $\varphi(x)$ 有界的假设可以改进为
\begin{align*}
|\varphi(x)|\leqslant Me^{Ax^2},\quad x\in\mathbb{R},
\end{align*}
这里 $A$ 和 $M$ 是正常数。此时 Poisson 公式 \eqref{def:poisson1d} 表示的函数 $u(x,t)$ 在区域 $\mathbb{R}\times(0,(4a^2A)^{-1})$ 上仍然是初值问题 \eqref{eq:heat-cauchy1d} 的 $C^\infty$ 解。
\end{remark}
\begin{proof}
当 $(x,t)\in\mathbb{R}_+^2$ 时，
\begin{align*}
|u(x,t)|&\leqslant \frac{1}{2a\sqrt{\pi t}}\int_{-\infty}^{+\infty}|\varphi(\xi)|e^{-\frac{(x-\xi)^2}{4a^2t}}\,\mathrm{d}\xi \\
&\leqslant \sup_{\mathbb{R}}|\varphi(\xi)|\frac{1}{\sqrt{\pi}}\int_{-\infty}^{+\infty}e^{-\eta^2}\,\mathrm{d}\eta \\
&\leqslant \sup_{\mathbb{R}}|\varphi(\xi)|.
\end{align*}
另外，固定任何一点 $(x_0,t_0)\in\mathbb{R}_+^2$，我们在其邻域 $(x_0-t_0,x_0+t_0)\times(\frac{t_0}{2},\frac{3t_0}{2})$ 内考虑 $u(x,t)$ 的微分性质。因此，$\frac{1}{\sqrt{t}}$ 在此邻域上无穷次可微，而 $e^{-\frac{(x-\xi)^2}{4a^2t}}$ 与任何次数的多项式 $P(\xi)$ 的乘积在 $\mathbb{R}$ 上可积，根据\hyperref[Basis of Analytics-theorem:含参量积分的可微性1]{含参量积分的可微性}我们知道 Poisson 公式 \eqref{def:poisson1d} 中的 $u(x,t)$ 关于 $x,t$ 求任意阶偏导数可以与关于 $\xi$ 求积分交换次序。由此我们得到 $u(x,t)$ 是 $C^\infty(\mathbb{R}_+^2)$ 中的函数，而且
\begin{align*}
\frac{\partial u(x,t)}{\partial t}-a^2\frac{\partial^2 u(x,t)}{\partial x^2}
=\int_{-\infty}^{+\infty}\left[\frac{\partial K(x-\xi,t)}{\partial t}-a^2\frac{\partial^2 K(x-\xi,t)}{\partial x^2}\right]\varphi(\xi)\,\mathrm{d}\xi=0.
\end{align*}
从而证明 $u(x,t)$ 满足热方程。

接着我们证明 $u(x,t)$ 满足初值条件。在 $u(x,t)$ 的积分表达式中作变量替换 $\eta=(\xi-x)/(2a\sqrt{t})$，则
\begin{align*}
u(x,t)=\frac{1}{\sqrt{\pi}}\int_{-\infty}^{+\infty}e^{-\eta^2}\varphi(x+2a\sqrt{t}\eta)\,\mathrm{d}\eta.
\end{align*}
利用 $\varphi$ 的有界性和连续性，上述积分在 $(x,t)\in\mathbb{R}\times\mathbb{R}_+$ 上关于 $x,t$ 一致收敛。因此
\begin{align*}
\lim_{\substack{x\to x_0\\ t\to 0^+}}u(x,t)
&=\frac{1}{\sqrt{\pi}}\int_{-\infty}^{+\infty}e^{-\eta^2}\varphi(x_0)\,\mathrm{d}\eta=\varphi(x_0).
\end{align*}
至此完成定理的证明.
\end{proof}

\begin{proposition}\label{proposition:一维热方程初值问题解的性质}
设 $\varphi(x)\in C(\mathbb{R})\cap L^\infty(\mathbb{R})$，$f(x,t)\equiv 0$,一维热方程初值问题的Poisson公式\eqref{def:poisson1d}表示的函数为
\begin{align*}
u(x,t)=\int_{-\infty}^{+\infty}K(x-\xi,t)\varphi(\xi)\,\mathrm{d}\xi
\end{align*}
\begin{enumerate}[(1)]
\item \textbf{奇偶性和周期性:}若 $\varphi$ 是奇（或偶，或周期为 $T$ 的）函数，$f\in C(\mathbb{R}\times(0,\infty))$ 有界,则$u(x,t)$
也是 $x$ 的奇（或偶，或周期为 $T$ 的）函数。

\item\label{proposition:一维热方程初值问题解的性质-2} \textbf{无限传播速度:}若存在$x_0\in \mathbb{R},\delta >0$,使得
\begin{gather*}
\varphi \left( x \right) >0,\quad \forall x\in \left( x_0-\delta ,x_0+\delta \right) ;
\\
\varphi \left( x \right) =0,\quad \forall x\in \mathbb{R} \backslash \left( x_0-\delta ,x_0+\delta \right) .
\end{gather*}
则对$\forall x\in \mathbb{R},t>0$,都有
\begin{align*}
u(x,t)=\int_{-\infty}^{+\infty}K(x-\xi,t)\varphi(\xi)\,\mathrm{d}\xi>0.
\end{align*}

\item \textbf{无穷次可微性:}无论$\varphi(x)$ 是否可微，当 $t>0$ 时，$u(x,t)$ 在 $\mathbb{R}\times\mathbb{R}_+$ 上必无穷次可微。
\end{enumerate}
\end{proposition}
\begin{remark}
性质\ref{proposition:一维热方程初值问题解的性质-2}的物理表述是：若无穷长杆的初始温度 $\varphi(x)$ 在小段 $(x_0-\delta,x_0+\delta)$ 大于零，且在其他地方 $\mathbb{R}\setminus(x_0-\delta,x_0+\delta)$ 等于零，则杆在 $t>0$ 时刻，在任何一点 $x$ 的温度
\begin{align*}
u(x,t)=\int_{-\infty}^{+\infty}K(x-\xi,t)\varphi(\xi)\,\mathrm{d}\xi>0.
\end{align*}
性质\ref{proposition:一维热方程初值问题解的性质-2}也说明，顷刻之间，杆的热量传递到杆上的任何一点，且离 $x_0$ 近的点，温度会升高多一些；而离 $x_0$ 远的点，温度会升高少一些。当然，这个结论与线性的 Fourier 热传导定律有关。
\end{remark}
\begin{proof}
\begin{enumerate}[(1)]
\item 

\item 

\item 事实上，当 $t>0$ 时，对于任意点 $x\in\mathbb{R}$，
\begin{align*}
\frac{\partial^{k+l}u}{\partial x^k\partial t^l}
=\int_{-\infty}^{+\infty}\frac{\partial^{k+l}}{\partial x^k\partial t^l}K(x-\xi,t)\varphi(\xi)\,\mathrm{d}\xi.
\end{align*}
这是因为对于任何正数 $\delta$，上述积分在区域 $\mathbb{R}\times(\delta,+\infty)$ 一致收敛，因而可以将求微商与求积分交换次序。
\end{enumerate}
\end{proof}


\begin{definition}
对 $x\in\mathbb{R}^n$，$t>0$，称函数
\begin{align*}
K(x,t)\triangleq 
\begin{cases}
\dfrac{1}{(4\pi a^2 t)^{n/2}}e^{-\frac{|x|^2}{4a^2t}}, & t>0,\\[0.4cm]
0, & t\leqslant 0
\end{cases}
\end{align*}
为 \textbf{$n$维热方程的Poisson 核}或\textbf{$n$维热方程的热核}，其中 $a>0$ 为常数.
\end{definition}

\begin{definition}
设 $\varphi\in C(\mathbb{R}^n)\cap L^\infty(\mathbb{R}^n)$，$f\in C(\mathbb{R}^n\times(0,\infty))$ 有界，$a>0$ 为常数. 对任意 $x\in\mathbb{R}^n$，$t>0$，定义
\begin{align}
u(x,t)\triangleq \int_{\mathbb{R}^n}K(x-\xi,t)\varphi(\xi)\,\mathrm{d}\xi
+\int_0^t\int_{\mathbb{R}^n}K(x-\xi,t-\tau)f(\xi,\tau)\,\mathrm{d}\xi\,\mathrm{d}\tau, \label{def:poisson-nd}
\end{align}
其中 $K$ 为 $n$ 维热方程的 Poisson 核. 称 \eqref{def:poisson-nd}式为 \textbf{$n$ 维热方程初值问题的Poisson 公式}.
\end{definition}

\begin{definition}
设 $\varGamma$ 是定义在 $\mathbb{R}^n\times(0,\infty)\times\mathbb{R}^n\times(0,\infty)$ 上的函数（其中 $t>\tau$），定义为
\begin{align*}
\varGamma(x,t;\xi,\tau)\triangleq K(x-\xi,t-\tau),\qquad t>\tau,
\end{align*}
其中 $K$ 为 $n$ 维热方程的 Poisson 核，即
\begin{align}
\varGamma(x,t;\xi,\tau)=
\begin{cases}
\dfrac{1}{(4\pi a^2(t-\tau))^{n/2}}\exp\left(-\dfrac{|x-\xi|^2}{4a^2(t-\tau)}\right), & t>\tau,\\[0.4cm]
0, & t\leqslant \tau.
\end{cases}
\end{align}
称 $\varGamma$ 为 \textbf{$n$ 维热方程的基本解}.
\end{definition}

\begin{theorem}
设 $\varphi(x)\in C(\mathbb{R}^n)\cap L^\infty(\mathbb{R}^n)$，$f(x,t)\in C^{0,1}(\mathbb{R}_+^{n+1})$，则$n$ 维热方程初值问题的Poisson 公式\eqref{def:poisson-nd}表示的函数
\begin{align*}
u(x,t)= \int_{\mathbb{R}^n}K(x-\xi,t)\varphi(\xi)\,\mathrm{d}\xi
+\int_0^t\int_{\mathbb{R}^n}K(x-\xi,t-\tau)f(\xi,\tau)\,\mathrm{d}\xi\,\mathrm{d}\tau
\end{align*}
是$n$ 维热方程初值问题\eqref{eq:多维热方程初值问题}的有界 $C^\infty(\mathbb{R}_+^{n+1})$ 解。
\end{theorem}










\end{document}